\section{Threats to Validity}
\label{sec:threats}
\paragraph{Construct validity.}
Runtime type errors and annotation-contract errors measure different
properties. An open-entry failure can describe an admissible symbolic
calling state without establishing a failure reachable from a repository
entry. Likewise, an annotation conflict need not cause a runtime exception
in an existing call. Our category-specific metrics, declared entries, and
cross-category links preserve these distinctions. Manual judgments about
reachability, intended interfaces, duplicate root causes, and exception
handling can affect precision. The source and behavioral evidence retained
with each report make these judgments inspectable; the adjudication
procedure is specified in Section~\ref{sec:setup}.

\paragraph{Internal validity.}
LLM-generated relations can omit legal type behaviors or introduce incorrect
ones. Conflict-driven expansion checks disputed relations, but unresolved
source semantics and incomplete context can still prevent a decision.
The method does not establish soundness or completeness for arbitrary
Python programs. The fixed pyflow call graph also bounds the callees
available for propagation and expansion; a missing target cannot be
recovered by revising a type-state relation alone. Dynamic attribute
behavior, reflection, native extensions, unknown dependency effects, and
bounded loop or recursive exploration can affect coverage. These cases
retain unresolved or unsupported outcomes rather than being treated as
confirmed absence of errors.

Ablation differences can arise from semantic coverage or budget consumption
as well as the representation being studied. Matched entries, rules, and
confirmation criteria help isolate these factors. We retain the causes of
undecided targets and compare intermediate failures in the relational and
direct-body alternatives. Differences between tools' dependency models and
supported Python features remain relevant to the whole-tool comparisons.

\paragraph{External validity.}
Public GitHub fixes emphasize errors that developers discovered and chose
to repair. Repository selection, issue wording, dependency availability,
and the required bug-fix pairing can therefore affect the sampled error
distribution. The recent collection interval reduces exposure to older
public examples but does not establish that the benchmark is absent from
model training data. Model training cutoffs and corpus membership may be
unknown. Independent discovery repositories test another setting, although
their domains and selected entries still bound generalization.

\paragraph{Conclusion validity.}
Each experimental configuration is run three times, and effectiveness
is reported using the best result. The selection metric, aggregation unit,
and tie-breaking rule are xxx. This best-of-three protocol reports selected
effectiveness rather than expected performance from a single run. It can favor methods with greater
run-to-run variability and depends on the selection metric and aggregation
unit. Consequently, the reported effectiveness cannot be interpreted as
mean reliability. Retaining all run records and accounting for all three
runs' cost exposes the associated variability and resource expenditure.
The selected run's cost is recorded separately.
Model revisions, decoding settings, and external service latency can also
change outcomes and cost.
